\section{Código Fuente}
\label{sec:codigo}

Esta sección presenta los fragmentos más relevantes del código fuente del proyecto, seleccionados para ilustrar las decisiones de implementación discutidas en secciones anteriores. Se omiten imports, docstrings extensos y código repetitivo para enfocarse en la lógica central.

% =============================================================================
% FRAGMENTOS DE model.py
% =============================================================================

\subsection{Definición del Modelo}

El constructor de la clase \texttt{SentimentRNN} define la arquitectura completa del modelo: capa de embedding, capas recurrentes apiladas (LSTM o GRU) y capa de clasificación lineal. La selección dinámica del tipo de RNN permite comparar ambas arquitecturas con una sola implementación.

% --- Clase SentimentRNN.__init__ — Definición de capas ---
\begin{minted}{python}
class SentimentRNN(nn.Module):
    def __init__(
        self,
        vocab_size: int,
        embedding_dim: int,
        hidden_size: int,
        num_layers: int,
        rnn_type: str,
        rnn_dropout: float,
        embed_dropout: float,
    ) -> None:
        super().__init__()

        if rnn_type not in ("LSTM", "GRU"):
            raise ValueError(
                f"Invalid rnn_type '{rnn_type}'. "
                f"Supported: 'LSTM', 'GRU'"
            )

        self.rnn_type = rnn_type
        self.hidden_size = hidden_size
        self.num_layers = num_layers

        # Embedding layer
        self.embedding = nn.Embedding(vocab_size, embedding_dim)
        self.embed_dropout = nn.Dropout(p=embed_dropout)

        # Stacked RNN (LSTM or GRU)
        rnn_cls = nn.LSTM if rnn_type == "LSTM" else nn.GRU
        self.rnn = rnn_cls(
            input_size=embedding_dim,
            hidden_size=hidden_size,
            num_layers=num_layers,
            dropout=rnn_dropout,
            batch_first=True,
        )

        # Output linear layer
        self.fc = nn.Linear(hidden_size, 1)
\end{minted}

La inicialización del bias de la compuerta de olvido (\textit{forget gate}) a 1.0 es una técnica propuesta por \cite{jozefowicz2015} que favorece la retención de información en las primeras etapas del entrenamiento, evitando que la celda LSTM descarte prematuramente el estado de memoria.

% --- Inicialización del bias de forget gate a 1.0 ---
\begin{minted}{python}
        # Initialize LSTM forget gate bias to 1.0 to encourage
        # remembering early in training (Jozefowicz et al., 2015).
        # The LSTM bias tensor is structured as:
        #   [input_gate | forget_gate | cell_gate | output_gate]
        # each segment of size hidden_size. We set the forget gate
        # segment (indices n//4 to n//2) to 1.0.
        if rnn_type == "LSTM":
            for name, param in self.rnn.named_parameters():
                if "bias" in name:
                    n = param.size(0)
                    param.data[n // 4 : n // 2].fill_(1.0)

        # Output linear layer
        self.fc = nn.Linear(hidden_size, 1)
\end{minted}

El método \texttt{forward} implementa el flujo de datos completo: la secuencia de índices pasa por la capa de embedding, luego por las capas recurrentes apiladas, y finalmente el estado oculto de la última capa se proyecta a un logit escalar para clasificación binaria.

% --- Método forward — Flujo embedding → RNN → fc ---
\begin{minted}{python}
    def forward(self, x: Tensor) -> Tensor:
        # Embedding: (batch, seq_len) -> (batch, seq_len, emb_dim)
        embedded = self.embedding(x)
        embedded = self.embed_dropout(embedded)

        # RNN forward pass
        # output: (batch, seq_len, hidden_size)
        _output, hidden = self.rnn(embedded)

        # Extract hidden state from the final layer
        if self.rnn_type == "LSTM":
            # h_n: (num_layers, batch, hidden_size)
            final_hidden = hidden[0][-1]
        else:
            final_hidden = hidden[-1]

        # Linear: (batch, hidden_size) -> (batch, 1)
        logits = self.fc(final_hidden)
        return logits
\end{minted}

% =============================================================================
% FRAGMENTOS DE train.py
% =============================================================================

\subsection{Entrenamiento}

La configuración del criterio de pérdida y el optimizador establece los componentes fundamentales del bucle de entrenamiento. Se utiliza \texttt{BCEWithLogitsLoss} que combina la función sigmoide con la entropía cruzada binaria en una sola operación numéricamente estable.

% --- Configuración de criterion y optimizer ---
\begin{minted}{python}
    # Move model to device
    model = model.to(device)

    # Loss function and optimizer
    criterion = nn.BCEWithLogitsLoss()
    optimizer = torch.optim.Adam(model.parameters(), lr=lr)

    # Early stopping state
    best_val_loss = float("inf")
    epochs_without_improvement = 0

    # Checkpoint path based on model's rnn_type
    checkpoint_path = os.path.join(
        checkpoint_dir, f"{model.rnn_type}_best.pt"
    )

    # Global step counter for gradient monitoring
    global_step = 0
\end{minted}

El bucle interno de entrenamiento ejecuta el paso forward, calcula la pérdida, propaga los gradientes y registra las normas L2 de los gradientes por capa durante los primeros 100 pasos globales. Este monitoreo permite detectar problemas de \textit{vanishing} o \textit{exploding gradients}.

% --- Bucle de entrenamiento: forward, backward, gradient monitoring ---
\begin{minted}{python}
        for inputs, labels in train_loader:
            inputs = inputs.to(device)
            labels = labels.to(device).float()

            # Forward pass
            optimizer.zero_grad()
            logits = model(inputs).squeeze(1)
            loss = criterion(logits, labels)

            # Backward pass
            loss.backward()

            # Record gradient norms for first 100 steps
            global_step += 1
            if global_step <= GRADIENT_MONITOR_STEPS:
                for name, param in model.named_parameters():
                    if param.grad is not None:
                        grad_norm = param.grad.data.norm(2).item()
                        if name not in history.gradient_norms:
                            history.gradient_norms[name] = []
                        history.gradient_norms[name].append(grad_norm)
\end{minted}

El \textit{gradient clipping} limita la norma L2 global de los gradientes a un valor máximo de 5.0, previniendo actualizaciones desproporcionadas que podrían desestabilizar el entrenamiento.

% --- Gradient clipping ---
\begin{minted}{python}
            # Gradient clipping: limit L2 norm of all gradients
            # to max_norm (5.0) to prevent exploding gradients.
            # Applied after loss.backward() and before
            # optimizer.step() to ensure stable updates.
            torch.nn.utils.clip_grad_norm_(
                model.parameters(), max_norm=clip_norm
            )

            # Optimizer step
            optimizer.step()

            # Track metrics
            running_loss += loss.item() * inputs.size(0)
            predictions = (torch.sigmoid(logits) > 0.5).long()
            correct += (predictions == labels.long()).sum().item()
            total += labels.size(0)
\end{minted}

El mecanismo de \textit{early stopping} compara la pérdida de validación de cada época con la mejor registrada. Si no hay mejora durante un número configurable de épocas (\texttt{patience}), el entrenamiento se detiene anticipadamente y se conserva el mejor checkpoint.

% --- Early stopping: comparación val_loss, guardado de checkpoint ---
\begin{minted}{python}
        # --- Early stopping check ---
        if epoch_val_loss < best_val_loss:
            best_val_loss = epoch_val_loss
            epochs_without_improvement = 0
            # Save best checkpoint
            torch.save(model.state_dict(), checkpoint_path)
            print(f"  Best model saved to {checkpoint_path}")
        else:
            epochs_without_improvement += 1
            print(
                f"  No improvement "
                f"({epochs_without_improvement}/{patience})"
            )
            if epochs_without_improvement >= patience:
                print(
                    f"  Early stopping after epoch {epoch + 1}"
                )
                break
\end{minted}

% =============================================================================
% FRAGMENTOS DE data.py
% =============================================================================

\subsection{Preprocesamiento de Datos}

La función \texttt{clean\_text} implementa el pipeline de limpieza de texto: elimina etiquetas HTML, convierte a minúsculas, remueve caracteres no alfabéticos y normaliza los espacios mediante tokenización y reconstrucción.

% --- Función clean_text — Preprocesamiento de texto ---
\begin{minted}{python}
def clean_text(text: str) -> str:
    """Clean review text: remove HTML, lowercase,
    remove punctuation, tokenize."""
    # Step 1: Remove HTML tags (e.g., <br />, <p>)
    text = re.sub(r'<[^>]+>', ' ', text)
    # Step 2: Convert to lowercase
    text = text.lower()
    # Step 3: Remove non-alpha characters
    # (keep only letters and spaces)
    text = re.sub(r'[^a-z\s]', '', text)
    # Step 4: Tokenize by whitespace and rejoin
    # to normalize multiple spaces
    tokens = text.split()
    return ' '.join(tokens)
\end{minted}

La función \texttt{pad\_sequence} garantiza que todas las secuencias tengan exactamente la misma longitud, truncando las que exceden el máximo y rellenando con ceros (token PAD) las más cortas. Esto es necesario para el procesamiento en lotes con tensores de dimensión fija.

% --- Función pad_sequence — Truncamiento y padding ---
\begin{minted}{python}
def pad_sequence(encoded: list[int], max_len: int) -> list[int]:
    """Pad or truncate a sequence to exactly max_len.

    If longer than max_len: truncate (keep first max_len tokens).
    If shorter than max_len: pad with zeros (PAD token = 0).
    This ensures all sequences in a batch have uniform length
    for efficient tensor operations.
    """
    if len(encoded) > max_len:
        return encoded[:max_len]
    elif len(encoded) < max_len:
        return encoded + [0] * (max_len - len(encoded))
    else:
        return list(encoded)
\end{minted}

% =============================================================================
% FRAGMENTOS DE evaluate.py
% =============================================================================

\subsection{Evaluación e Inferencia}

El proceso de inferencia carga el mejor checkpoint guardado durante el entrenamiento, ejecuta el modelo en modo evaluación sin cálculo de gradientes, y aplica un umbral de 0.5 sobre la salida sigmoide para obtener las predicciones binarias finales.

% --- Inferencia: carga de checkpoint, forward sin gradientes, sigmoid > 0.5 ---
\begin{minted}{python}
def evaluate_model(
    model: nn.Module,
    test_loader: DataLoader,
    device: torch.device,
    checkpoint_path: str,
    rnn_type: str,
) -> EvaluationResults:
    # Load checkpoint weights
    state_dict = torch.load(
        checkpoint_path, map_location=device, weights_only=True
    )
    model.load_state_dict(state_dict)

    # Move model to device and set eval mode
    model = model.to(device)
    model.eval()

    all_predictions: list[int] = []
    all_labels: list[int] = []

    # Inference with no gradient computation
    with torch.no_grad():
        for inputs, labels in test_loader:
            inputs = inputs.to(device)
            labels = labels.to(device)

            logits = model(inputs).squeeze(1)

            # Apply sigmoid and threshold at 0.5
            probs = torch.sigmoid(logits)
            preds = (probs > 0.5).long()

            all_predictions.extend(preds.cpu().numpy().tolist())
            all_labels.extend(labels.cpu().numpy().tolist())
\end{minted}
